% ============ Appendix B. Some homological algebra (原书 pp.274-282) ============
\chapter*{Appendix B}
\addcontentsline{toc}{chapter}{Appendix B: Some homological algebra}
\markboth{Appendix B}{Some homological algebra}
\renewcommand{\thetheorem}{B\arabic{theorem}}
\setcounter{theorem}{0}

\vspace*{-0.5cm}
\begin{center}
  {\Large\bfseries Some homological algebra}
\end{center}
\vspace{0.3cm}

Let $A$ be a ring; by a map from an $A$-module into another we mean an
$A$-linear map.

\medskip
\noindent\textbf{Complexes}

By a complex we mean a sequence
\[
  \cdots \to K_{n+1} \xrightarrow{d_{n+1}} K_n \xrightarrow{d_n} K_{n-1} \to \cdots
\]
of $A$-modules and $A$-linear maps such that $d_{n-1} \circ d_n = 0$ for every $n$.
This complex is written $K_\bullet$. Since $\Im(d_{n+1}) \subset \Ker(d_n)$ we can define
$H_n(K_\bullet) = \Ker(d_n)/\Im(d_{n+1})$ to be the \emph{homology} of $K_\bullet$ in dimension $n$. To
say that $H_n(K_\bullet) = 0$ for all $n$ is to say that $K_\bullet$ is exact. We also consider
complexes in which the indices go the other way, $\cdots \to K^n \to K^{n+1} \to \cdots$,
and for these we write $K^\bullet$ for the complex, and
$H^n(K^\bullet) = \Ker(d_n)/\Im(d_{n-1})$ for the \emph{cohomology} of $K^\bullet$ in dimension $n$.
From now on we often omit the indices, writing $d$ for $d_n$. We call $d$ the
\emph{differential} of the complex $K_\bullet$.

A morphism $f\colon K_\bullet \to K'_\bullet$ of complexes is a family $f = (f_n)_{n \in \mathbb{Z}}$ of
$A$-linear maps $f_n\colon K_n \to K'_n$ satisfying $d' \circ f_n = f_{n-1} \circ d$, or in other
words a commutative diagram
\[
\begin{tikzcd}
\cdots \arrow[r] & K_n \arrow[r] \arrow[d, "f_n"] & K_{n-1} \arrow[r] \arrow[d, "f_{n-1}"] & K_{n-2} \arrow[r] \arrow[d, "f_{n-2}"] & \cdots \\
\cdots \arrow[r] & K'_n \arrow[r] & K'_{n-1} \arrow[r] & K'_{n-2} \arrow[r] & \cdots
\end{tikzcd}
\]
In an obvious way, $f$ induces a linear map $H_n(K_\bullet) \to H_n(K'_\bullet)$ between the
homology modules in each dimension; this is often written $f_{n*}$, or simply $f$
if there is no fear of confusion. If $f, g\colon K_\bullet \to K'_\bullet$ are two morphisms we say
that $f$ and $g$ are \emph{homotopic} (denoted $f \sim g$) if for each $n$ there is a linear map
$h_n\colon K_n \to K'_{n+1}$ such that
\[
  f_n - g_n = d'h_n + h_{n-1}d.
\]
If this happens then $f$ and $g$ induce the same map $H_n(K_\bullet) \to H_n(K'_\bullet)$ on
homology. Two complexes $K_\bullet$ and $K'_\bullet$ are said to be \emph{homotopy equivalent} if
there exist morphisms $f\colon K_\bullet \to K'_\bullet$ and $g\colon K'_\bullet \to K_\bullet$ such that $gf \sim 1$,
and $fg \sim 1_{K'}$, where $1$ denotes the identity map $K_\bullet \to K_\bullet$. Homotopy
equivalent complexes have the same homology.

A sequence of complexes
\[
  0 \to K'_\bullet \xrightarrow{f} K_\bullet \xrightarrow{g} K''_\bullet \to 0
\]
is said to be \emph{exact} if
\[
  0 \to K'_n \xrightarrow{f_n} K_n \xrightarrow{g_n} K''_n \to 0
\]
is exact for every $n$. In this case, a connecting homomorphism
$\delta_n\colon H_n(K''_\bullet) \to H_{n-1}(K'_\bullet)$ is defined as follows: for $\zeta \in H_n(K''_\bullet)$, choose
$x \in \Ker d''_n$ representing $\zeta$, and take $y \in K_n$ such that $g(y) = x$; then since
$g(dy) = 0$ there is a well-determined $z \in K'_{n-1}$ for which $f_{n-1}(z) = dy$, and
$dz = 0$. The class $\xi \in H_{n-1}(K')$ represented by $z$ can easily be seen to depend
only on $\zeta$, and $\delta_n$ is defined by $\delta_n(\zeta) = \xi$. The following sequence is then
exact:
\[
  \cdots \to H_n(K'_\bullet) \xrightarrow{f} H_n(K_\bullet) \xrightarrow{g} H_n(K''_\bullet)
      \xrightarrow{\delta_n} H_{n-1}(K'_\bullet) \to \cdots
\]
The proof does not require anything new, and is well-known, so that we omit
it; this should be thought of as a fundamental theorem of homology theory.
The above sequence is called the \emph{homology long exact sequence} of the short exact
sequence $0 \to K'_\bullet \to K_\bullet \to K''_\bullet \to 0$.

\medskip
\noindent\textbf{Double complexes}

A double complex of $A$-modules is a doubly indexed family
$K_{\bullet\bullet} = (K_{pq})_{p, q \in \mathbb{Z}}$ of $A$-modules, with two sets of $A$-linear maps
$d'\colon K_{pq} \to K_{p-1,q}$ and $d''\colon K_{pq} \to K_{p,q-1}$ for which $d'd' = 0$, $d''d'' = 0$ and
$d'd'' = d''d'$. Given a double complex $K_{\bullet\bullet}$, if we set
\[
  K_n = \bigoplus_{p + q = n} K_{pq}
\]
and define $d_n\colon K_n \to K_{n-1}$ by
\[
  dx = d'x + (-1)^p d''x \quad \text{if } x \in K_{pq},
\]
then since $dd = 0$, the $\{K_n\}$ form an ordinary complex with differential $d$. The
homology of this complex is called the homology of $K_{\bullet\bullet}$, and written
$H_\bullet(K_{\bullet\bullet})$, or simply $H(K)$.

To treat homology and cohomology in a unified manner, we fix the following
convention on raising and lowering indices: $K_{pq} = K^{-p,-q} = K^{p,q} = K^{-p,-q}$.
For example, given a double complex $(K^{pq})$ with
$d'\colon K^{pq} \to K^{p+1,q}$ and $d''\colon K^{pq} \to K^{p,q+1}$, we think of $K^{pq}$ as
$K_{-p,-q}$, and set
\[
  K^n = \bigoplus_{p + q = n} K^{pq}.
\]

The basic technique for studying the homology of double complexes is spectral
sequences, but we leave this to specialist texts, and only consider here the
extreme cases which we will use later.

We can fix the first index $p$ in $K_{\bullet\bullet}$ getting a complex $K_{p\bullet}$:
\[
  \cdots \to K_{p,q+1} \xrightarrow{d''} K_{pq} \xrightarrow{d''} K_{p,q-1} \to \cdots
\]
Similarly, for fixed $q$, $d'$ defines a complex $K_{\bullet q}$.

Now suppose that $K_{\bullet\bullet}$ satisfies the condition $K_{pq} = 0$ if $p$ or $q < 0$ (a
first quadrant double complex). We set $H_0(K_{p\bullet}) = K_{p0}/d''K_{p1} = X_p$; then
$d'$ induces a map $X_p \to X_{p-1}$, making the $X_p$ into a complex $X_\bullet$. Similarly,
$d''$ makes the $H_q(K_{\bullet q}) = Y_q$ into a complex $Y_\bullet$. In this notation we have
the following theorem.

\begin{theorem}\label{thm:B1}
Suppose that the double complex $K_{\bullet\bullet}$ satisfies the conditions
\[
  K_{pq} = 0 \quad \text{for } p \text{ or } q < 0,
\]
and
\[
  H_q(K_{p\bullet}) = 0 \quad \text{for } q > 0 \text{ and all } p.
\]
Then in the above notation we have
\[
  H_n(K) \cong H_n(X_\bullet) \quad \text{for all } n.
\]
If in addition we have $H_p(K_{\bullet q}) = 0$ for $p > 0$ and all $q$ then
\[
  H_n(X_\bullet) = H_n(K) = H_n(Y_\bullet).
\]
\end{theorem}

\begin{proof}[Sketch proof]
Write $a_{ij}$ to denote an element of $K_{ij}$. We define a map
$\Phi\colon K_n \to X_n$ by taking $a = \sum_{i=0}^{n} a_{n-i,i} \in K_n$ into $\varphi(a_{n0}) \in X_n$, where
$\varphi\colon K_{n0} \to X_n$ denotes the canonical map. Then $\Phi$ is a morphism of
complexes, and we prove that it induces an isomorphism on homology.

Let $x \in X_n$ with $d'x = 0$. We can take $a_{n0}$ such that $x = \varphi(a_{n0})$, and
then since $\varphi(d'a_{n0}) = d'x = 0$ there exists $a_{n-1,1}$ such that
$d'a_{n0} = d''a_{n-1,1}$. In turn, since
$d''(d'a_{n-1,1}) = d'(d''a_{n-1,1}) = d'd'a_{n0} = 0$ and since $H_1(K_{n-2,\bullet}) = 0$ there
exists $a_{n-2,2}$ such that $da_{n-1,1} = d''a_{n-2,2}$; then proceeding as before, we can
choose $a_{n-i,i}$ for $0 \le i \le n$ such that $d'a_{n-i,i} = d''a_{n-i-1,i+1}$ for
$0 \le i < n$. Then for a suitable choice of signs,
$a = \sum_i a_{n-i,i} \in K_n$ satisfies $da = 0$ and $\Phi(a) = x$, and this proves that
$\Phi$ induces a surjection $H_n(K) \to H_n(X_\bullet)$. The proof that this is also
injective is similar. The second part follows by symmetry from the first.
\end{proof}

Dually, we have the following theorem for cohomology.

\begin{theorem}\label{thm:B2}
Suppose that the double complex $K^{\bullet\bullet}$ satisfies
\[
  K^{pq} = 0 \quad \text{for } p \text{ or } q < 0
\]
and
\[
  H^q(K^{p\bullet}) = 0 \quad \text{for } q > 0 \text{ and all } p.
\]
Then making $X^p = \Ker(d''\colon K^{p0} \to K^{p1})$ into a complex $X^\bullet$ by means of
$d'$, we have
\[
  H^n(K) \cong H^n(X^\bullet),
\]
If in addition $H^p(K^{\bullet q}) = 0$ for $p > 0$ and all $q$ then
\[
  H^n(Y^\bullet) \cong H^n(K) \cong H^n(X^\bullet),
\]
where $Y^\bullet$ is the complex made from $Y^q = \Ker(d'\colon K^{0q} \to K^{1q})$.
\end{theorem}

\begin{proof}
We leave the proof as a suitable exercise for the reader.
\end{proof}

\medskip
\noindent\textbf{Projective and injective modules}

An $A$-module $P$ is said to be \emph{projective} if it satisfies the following condition:
for any surjection $f\colon M \to N$ and any map $g\colon P \to N$, there exists a lifting
$h\colon P \to M$ such that $g = fh$. A free module is projective, and we can
characterise projective modules as direct summands of free modules. Indeed, if
we express a projective module $P$ as a quotient $P = F/G$ of a free module $F$
then the identity map $P \to P$ has a lifting such that $P \to F \to P$ is the identity
map, and then $F \cong P \oplus G$. Reversing the arrows and replacing surjection by
injection in the definition of projective module, we get the definition of
\emph{injective} module. There is no dual notion to that of a free module, so that
injective modules do not have any very simple characterisation, but we can
easily prove the following theorem using Zorn's lemma.

\begin{theorem}\label{thm:B3}
A necessary and sufficient condition for an $A$-module $I$ to be injective is that
for any ideal $\mathfrak{a}$ of $A$, and any map $\varphi\colon \mathfrak{a} \to I$, it is possible to extend
$\varphi$ to a map from the whole of $A$ to $I$.
\end{theorem}

Any $A$-module can be written as a quotient of a projective module (take for
example a free module). Dually to this, any module can be embedded into an
injective module; the proof of this is a little tricky, and we leave it to more
specialist textbooks. Given a module $M$, consider a surjection
$P_0 \to M$ from a projective module $P_0$ to $M$; letting $K_1$ be the kernel, we get an
exact sequence $0 \to K_1 \to P_0 \to M \to 0$. In the same way, we construct for $K_1$
an exact sequence $0 \to K_2 \to P_1 \to K_1 \to 0$ with $P_1$ projective, and proceeding as
before we get exact sequences $0 \to K_i \to P_{i-1} \to K_{i-1} \to 0$ for $i = 1, 2, \ldots$ with
$P_i$ projective. The resulting complex
\[
  P_\bullet\colon \cdots \to P_n \to P_{n-1} \to \cdots \to P_1 \to P_0 \to 0
\]
is called a \emph{projective resolution} of $M$. Since by construction this becomes an exact
sequence on replacing the final $P_0 \to 0$ by $P_0 \to M \to 0$, we have $H_n(P_\bullet) = 0$
for $n > 0$ and $H_0(P_\bullet) = M$. In the case that $A$ is Noetherian and $M$ is finite,
we can take $P_0$ to be a free module of finite rank, and then $K_1$ is again
finitely generated. Proceeding in the same way, we see that $M$ has a projective
resolution in which each $P_n$ is a finite free module.

Dually, for any $A$-module $M$ there exists an exact sequence of the form
$0 \to M \to Q^0 \to Q^1 \to \cdots$ with each $Q^n$ an injective module. The complex
$Q^\bullet\colon 0 \to Q^0 \to Q^1 \to \cdots$ is called an \emph{injective resolution} of $M$; it
satisfies $H^0(Q^\bullet) = M$ and $H^n(Q^\bullet) = 0$ for $n > 0$.

If $f\colon M \to N$ is a map of $A$-modules and $P_\bullet$, $P'_\bullet$ are projective
resolutions of $M$ and $N$ then there exists a morphism of complexes
$\varphi\colon P_\bullet \to P'_\bullet$ for which $f\varepsilon = \varepsilon'\varphi_0$, that is, a commutative diagram
\[
\begin{tikzcd}
\cdots \arrow[r] & P_1 \arrow[r] \arrow[d, "\varphi_1"] & P_0 \arrow[r] \arrow[d, "\varphi_0"] & M \arrow[r] \arrow[d, "f"] & 0 \\
\cdots \arrow[r] & P'_1 \arrow[r] & P'_0 \arrow[r] & N \arrow[r] & 0
\end{tikzcd}
\]
The existence of $\varphi_0, \varphi_1, \ldots$ can easily be proved successively, using the fact
that the $P_n$ are projective and the exactness of the lower sequence. Up to
homotopy, this $\varphi$ is unique, that is if $\varphi$ and $\psi$ both have the given property
then $\varphi \sim \psi$. We leave the proof of this to the reader. From this it follows
that any two projective resolutions of $M$ are homotopy equivalent. Exactly the
same thing holds for injective resolutions.

\medskip
\noindent\textbf{The Tor functors}

Let $M$ and $N$ be $A$-modules and $P_\bullet$, $Q_\bullet$ projective resolutions of $M$ and
$N$, respectively. We write $P_\bullet \otimes N$ for the complex obtained by tensoring
$P_\bullet$ through with $N$:
\[
  P_\bullet \otimes N\colon \cdots \to P_1 \otimes N \to P_0 \otimes N \to 0.
\]
The complex $M \otimes Q_\bullet$ is constructed similarly. Moreover, we can define a
double complex $K_{\bullet\bullet}$ by $K_{pq} = P_p \otimes_A Q_q$, with the obvious definitions of
$d'$, $d''$. Each $P_p$ is a direct summand of a free module, and is therefore flat
(that is performing $\otimes P_p$ takes exact sequences into exact sequences). Thus
$H_q(K_{p\bullet}) = H_q(P_p \otimes Q_\bullet) = 0$ for $q > 0$, and
$H_0(K_{p\bullet}) = H_0(P_p \otimes Q_\bullet) = P_p \otimes N$. In exactly the same way,
$H_p(K_{\bullet q}) = 0$ for $p > 0$ and $H_0(K_{\bullet q}) = M \otimes Q_q$, and therefore by
Theorem~\hyperref[thm:B1]{B1} $H_n(P_\bullet \otimes N) \cong H_n(K_{\bullet\bullet}) \cong H_n(M \otimes Q_\bullet)$. This
module (defined uniquely up to isomorphism) is written $\Tor^A_n(M, N)$; it is
independent of the choice of the projective resolutions of $M$ and $N$ chosen,
since if $P_\bullet$ and $P'_\bullet$ are two projective resolutions, we have $P_\bullet \sim P'_\bullet$, and
therefore $P_\bullet \otimes N \sim P'_\bullet \otimes N$.

The Tor functors have the following properties, (all of which can be proved
directly from the definition):
\begin{enumerate}[(1)]
  \item $\Tor^A_0(M, N) = M \otimes_A N$;
  \item if $M$ is flat then $\Tor^A_n(M, N) = 0$ for any $N$ and $n > 0$;
  \item $\Tor^A_n(M, N) \cong \Tor^A_n(N, M)$;
  \item $\Tor^A_n(M, N)$ is a covariant functor in both of its entries, and each short
        exact sequence $0 \to M' \to M \to M'' \to 0$ leads to a long exact sequence
        \[
          \cdots \to \Tor^A_n(M', N) \to \Tor^A_n(M, N) \to \Tor^A_n(M'', N)
          \to \Tor^A_{n-1}(M', N) \to \cdots \to \Tor^A_0(M'', N)
          \to M' \otimes N \to M \otimes N \to M'' \otimes N \to 0.
        \]
  \item If $(N_\lambda, f_{\mu\lambda})$ is a direct system of $A$-modules then
        \[
          \Tor^A_n(M, \varinjlim N_\lambda) = \varinjlim \Tor^A_n(M, N_\lambda).
        \]
\end{enumerate}

\medskip
\noindent\textbf{The Ext functors}

Let $M$ and $N$ be $A$-modules. The functor $\Hom_A(M, -)$ is left-exact, that is it
takes an exact sequence $0 \to N' \to N \to N'' \to 0$ into an exact sequence
$0 \to \Hom_A(M, N') \to \Hom_A(M, N) \to \Hom_A(M, N'')$; and $M$ is projective if and
only if $\Hom_A(M, -)$ is exact. In addition, $\Hom_A(-, N)$ is left-exact, in the sense
that it takes an exact sequence $0 \to M' \to M \to M'' \to 0$ into an exact sequence
$0 \to \Hom_A(M'', N) \to \Hom_A(M, N) \to \Hom_A(M', N)$, and $N$ is injective if and
only if $\Hom_A(-, N)$ is exact.

Choose a projective resolution $P_\bullet$ of $M$ and an injective resolution $Q_\bullet$ of $N$;
we define a double complex $K^{\bullet\bullet}$ by $K^{pq} = \Hom_A(P_p, Q^q)$, and construct
the two complexes
\[
  \Hom_A(M, Q^\bullet)\colon 0 \to \Hom_A(M, Q^0) \to \Hom_A(M, Q^1) \to \cdots
\]
and
\[
  \Hom_A(P_\bullet, N)\colon 0 \to \Hom_A(P_0, N) \to \Hom_A(P_1, N) \to \cdots.
\]
Then by Theorem~\hyperref[thm:B2]{B2} we get
\[
  H^n(\Hom_A(M, Q^\bullet)) \cong H^n(K^{\bullet\bullet}) \cong H^n(\Hom_A(P_\bullet, N)).
\]
Identifying these three, we write $\Ext^n_A(M, N)$. As with Tor, this does not depend
on the choice of $P_\bullet$ and $Q^\bullet$.

The main properties of the Ext functors are as follows:
\begin{enumerate}[(1)]
  \item $\Ext^0_A(M, N) = \Hom_A(M, N)$,
  \item If $M$ is projective, or if $N$ is injective, then $\Ext^n_A(M, N) = 0$ for $n > 0$;
  \item $\Ext^n_A(M, N)$ is a contravariant functor in $M$ and a covariant functor
        in $N$. A short exact sequence $0 \to M' \to M \to M'' \to 0$ gives rise to
        a long exact sequence
        \[
          0 \to \Hom_A(M'', N) \to \Hom_A(M, N) \to \Hom_A(M', N)
          \to \Ext^1_A(M'', N) \to \Ext^1_A(M, N) \to \Ext^1_A(M', N)
          \to \Ext^2_A(M'', N) \to \cdots,
        \]
        and a short exact sequence $0 \to N' \to N \to N'' \to 0$ gives rise to a long
        exact sequence
        \[
          0 \to \Hom_A(M, N') \to \Hom_A(M, N) \to \Hom_A(M, N'')
          \to \Ext^1_A(M, N') \to \Ext^1_A(M, N) \to \Ext^1_A(M, N'')
          \to \Ext^2_A(M, N') \to \cdots;
        \]
  \item $M$ is projective $\iff \Ext^1_A(M, N) = 0$ for all $N$,
        and
        \[
          N \text{ is injective } \iff \Ext^1_A(M, N) = 0 \text{ for all } M.
        \]
\end{enumerate}

\medskip
\noindent\textbf{Projective and injective dimensions}

If $M$ is an $A$-module for which there exists a projective resolution $P_\bullet$ with
$P_n = 0$ for $n > d$, but such that $P_d \ne 0$ for any choice of projective resolution
$P_\bullet$, then we say that $M$ has \emph{projective dimension} $d$, and write
$\projdim M = d$. If there is no such $d$ then we write $\projdim M = \infty$. The
\emph{injective dimension} $\injdim M$ is defined in the same way using injective
resolutions. Clearly $\projdim M = 0$ if and only if $M$ is projective, and
$\injdim M = 0$ if and only if $M$ is injective.

For a projective resolution $P_\bullet$ of $M$ and some $i > 0$, let $K_i$ denote the image of
$P_i \to P_{i-1}$; then $\cdots \to P_n \to P_{n-1} \to \cdots \to P_i \to 0$ is a projective
resolution of $K_i$, so that for $n > i$ we have $\Ext^n_A(M, N) = \Ext^{n-i}_A(K_i, N)$.
Now if $\Ext^{d+1}_A(M, N) = 0$ for all $N$, we have $\Ext^1_A(K_d, N) = 0$ for all $N$, and
hence $K_d$ is projective, so that $0 \to K_d \to P_{d-1} \to \cdots \to P_0 \to 0$ is also
a projective resolution of $P_0$, and $\projdim M \le d$. Conversely, if
$\projdim M \le d$ then obviously $\Ext^n_A(M, N) = 0$ for $n > d$.

Similarly, $\injdim N \le d \iff \Ext^{d+1}_A(M, N) = 0$ for all $M$.

\medskip
\noindent\textbf{Derived functors}

As we have just seen, the definition of functors like Tor and Ext can be given
using a resolution of just one entry. For instance, write $T$ for the functor
$\Hom_A(-, N)$, and let $P_\bullet$ be a projective resolution of a given module $M$;
construct the complex $T(P_\bullet)\colon \cdots \to T(P_n) \to T(P_{n-1}) \to \cdots \to T(P_0) \to 0$,
and take the cohomology $H^n(T(P_\bullet))$. Setting $R^nT(M) = H^n(T(P_\bullet))$ defines
a functor $R^nT$ in $M$, which we call the \emph{right derived functor} of the left exact
functor $T$. In the present case we have $R^nT = \Ext^n_A(-, N)$, but we can in
general define the right derived functor of a left exact contravariant functor
in the same way.

The right derived functor is uniquely determined by the following three
properties: (1) $R^0T = T$, (2) if $M$ is projective then $R^nT(M) = 0$ for all $n > 0$,
and (3) a short exact sequence $0 \to M' \to M \to M'' \to 0$ gives rise to a `natural'
long exact sequence
\[
  0 \to T(M'') \to T(M) \to T(M')
  \to R^1T(M'') \to R^1T(M) \to R^1T(M')
  \to R^2T(M'') \to \cdots
\]
(For the meaning of `natural' see a textbook on homological algebra.) For a
left exact covariant functor, we have to replace `projective' by `injective' in
the above. For right exact functors we can define a \emph{left derived functor} by
taking a projective resolution in the covariant case and an injective resolution
in the contravariant case.

\medskip
\noindent\textbf{Injective hull}

Let $L$ be an $A$-module and $M \subset L$ a submodule; we say that $L$ is an
\emph{essential extension} of $M$ if $N \cap M \ne 0$ for every non-zero submodule $N \subset L$,
or equivalently if
\[
  0 \ne x \in L \Rightarrow \text{there exists } a \in A \text{ such that } 0 \ne ax \in M.
\]

\begin{theorem}\label{thm:B4}
An $A$-module $M$ is injective if and only if it has no essential extensions
except $M$ itself.
\end{theorem}

We leave the proof to the reader. Now suppose that $M$ is a given $A$-module,
and choose an injective module $I$ with $M \subset I$. If we let $E$ be a maximal
element among all essential extensions of $M$ in $I$ then by the above theorem
$E$ is injective. An injective module $E$ such that $M \subset E$ is an essential
extension is called an \emph{injective hull} of $M$, and written $E(M)$ or $E_A(M)$; this
notion plays an important role in \S\ref{sec:18}. If $E$ and $E'$ are injective hulls of
$M$ then it is easy to see that there is an isomorphism $\varphi\colon E \to E'$ which
fixes the elements of $M$, although $\varphi$ itself is not necessarily unique.

Let $M$ be an $A$-module. Take an injective hull $I^0$ of $M$, and set
$K^1 = I^0/M$. Take an injective hull $I^1$ of $K^1$, and set $K^2 = I^1/K^1$.
Proceeding in the same way we obtain an injective resolution
$0 \to I^0 \to I^1 \to I^2 \to \cdots$ of $M$, which is called a \emph{minimal injective
resolution} of $M$.

The following two propositions are both famous and useful; the proofs are easy.

\medskip
\noindent\textbf{The five lemma.} Let
\[
\begin{tikzcd}
A \arrow[r] \arrow[d, "f_1"] & B \arrow[r] \arrow[d, "f_2"] & C \arrow[r] \arrow[d, "f_3"] & D \arrow[r] \arrow[d, "f_4"] & E \arrow[d, "f_5"] \\
A' \arrow[r] & B' \arrow[r] & C' \arrow[r] & D' \arrow[r] & E'
\end{tikzcd}
\]
be a commutative diagram of modules with exact rows. Then
\begin{enumerate}[(1)]
  \item $f_2$ surjective, and $f_1$ and $f_4$ injective $\Rightarrow f_3$ is injective;
  \item $f_4$ injective, and $f_2$ and $f_5$ surjective $\Rightarrow f_3$ is surjective.
\end{enumerate}

\medskip
\noindent\textbf{The snake lemma.} Let
\[
\begin{tikzcd}
A \arrow[r] \arrow[d, "\alpha"] & B \arrow[r] \arrow[d, "\beta"] & C \arrow[r] \arrow[d, "\gamma"] & 0 \\
0 \arrow[r] & A' \arrow[r] & B' \arrow[r] & C'
\end{tikzcd}
\]
be a commutative diagram of modules with exact rows. Then there exists an
exact sequence of the form
\[
  \Ker(\alpha) \to \Ker(\beta) \to \Ker(\gamma) \xrightarrow{\delta}
  \Coker(\alpha) \to \Coker(\beta) \to \Coker(\gamma).
\]

\medskip
\noindent\textbf{Tensor product of complexes}

Given two complexes of $A$-modules $K_\bullet$ and $L_\bullet$ the tensor product
$K_\bullet \otimes_A L_\bullet$ is defined as follows: firstly, set
\[
  (K \otimes L)_n = \bigoplus_{p + q = n} K_p \otimes_A L_q,
\]
and define the differential $d$ by setting
\[
  d(x \otimes y) = dx \otimes y + (-1)^p x \otimes dy
\]
for $x \in K_p$ and $y \in L_q$. In other words, $K \otimes L$ is the (single) complex obtained
from the double complex $W_{\bullet\bullet}$, where $W_{pq} = K_p \otimes L_q$.

There is an isomorphism of complexes $K \otimes L \cong L \otimes K$ obtained by sending
$x \otimes y$ into $(-1)^{pq}y \otimes x$ for $x \otimes y \in K_p \otimes L_q$. For a third complex of
$A$-modules $M$, the associative law holds:
\[
  (K \otimes L) \otimes M = K \otimes (L \otimes M).
\]
Hence the tensor product $K^{(1)} \otimes \cdots \otimes K^{(r)}$ of a finite number of complexes
can be defined by induction. This is used in \S\ref{sec:16}.

The information on homological algebra given above should be adequate for
the purpose of reading this book. However, a student intending to become a
specialist in algebra or geometry will require rather more detailed knowledge,
including the theory of spectral sequences. We mention here just three
representative references, two books by the originators of homological algebra
and category theory:
\begin{quote}
  H.~Cartan and S.~Eilenberg, \emph{Homological Algebra}, Princeton, 1956,\\
  S.~MacLane, \emph{Homology}, Springer, 1963,
\end{quote}
together with A.~Grothendieck's paper
\begin{quote}
  \emph{Sur quelques points d'alg\`ebre homologique}, T\^ohoku Math.~J.~9 (1957), 119--221.
\end{quote}
